\section{Competitive Dynamics}
\label{sec:dynamics}
Section~\ref{sec:monoculture} characterized when developers' equilibrium capability profiles exhibit monoculture or specialization. 
We now turn from the direction of capabilities at equilibrium to the evolution of performance levels over time. 
In particular, we ask when competing developers become closer in performance as investment accumulates, even if differences in their capability profiles persist.

Let $\mathbf{X}_{i,t} = \sum_{s=1}^t \mathbf{x}_{i,s}$
denote developer $i$'s cumulative investment through time $t$.
Additionally, let $\boldsymbol{\theta}_i^t = F_i(\mathbf X_i^t)$ denote developer $i$'s capability vector at time $t$.
At each time step $t$, each developer chooses $\mathbf{x}_{i,t}$ to maximize its current-period utility:
\[
\frac{1}{N-1}
\sum_{i'\neq i}\sum_k
w_k
\sigma\!\left(
F_i(\mathbf{X}_{i,t-1}+\mathbf{x}_{i,t})_k
-
F_{i'}(\mathbf{X}_{i',t-1}+\mathbf{x}_{i',t})_k
\right)
-
c(\|\mathbf{x}_{i,t}\|_1).
\]
This abstraction isolates how repeated investment changes relative performance without introducing intertemporal strategic considerations. 
We ask whether developers’ relative performance converges over time.
We do not expect competition to guarantee convergence: developers may remain separated because their cumulative investments diverge or because differences in their production functions persist at scale.
Rather than characterize the full dynamic equilibrium, we isolate conditions on investment trajectories and production functions under which relative performance must converge.

\begin{restatable}[Catch-Up and Long-Run Crowding]
{proposition}{AsymptoticBunching}
\label{prop:AsymptoticBunching}

Let $\boldsymbol{\theta}_i^t = F_i(\mathbf{X}_i^t)$ denote developer $i$'s capability vector at time $t$, and define its category-specific gap to the frontier as $g_{i,k}^t=\max_j \theta_{j,k}^t-\theta_{i,k}^t.$
Let $\mathbf g^t\in\mathbb R_+^{NK}$ stack these gaps across developers and
categories at timestep $t$, and let
\[
\Delta\boldsymbol{\theta}^t
=
\left(
\boldsymbol{\theta}_1^{t+1}-\boldsymbol{\theta}_1^t,
\ldots,
\boldsymbol{\theta}_N^{t+1}-\boldsymbol{\theta}_N^t
\right)
\in\mathbb R^{NK}.
\]

Suppose each $F_i$ is continuously differentiable with
\[
\max_{\ell\in[K]} \left\| \frac{\partial F_i(\mathbf x)}{\partial x_\ell} \right\|_2 \to 0 \qquad \text{as }\|\mathbf x\|_1\to\infty, \quad \forall i.
\]
Suppose there exists $\beta>0$ such that, for sufficiently large $t$, $\left\langle
\mathbf g^t,\Delta\boldsymbol{\theta}^t
\right\rangle
\ge
\beta\|\mathbf g^t\|_2^2$.
Then we have that $\|\mathbf g^t\|_2\to 0$, and $\max_{i'}\theta_{{i'},k}^t-\theta_{i,k}^t
\to 0$ for all developers $i$ and categories $k$. 
\end{restatable}

The competitive dynamics generate a sequence of relative performance gaps and per-period investments that decrease over time, while the catch-up condition requires improvements to be directed toward those gaps (proof in Appendix~\ref{app:AsymptoticBunching}).

We next give one set of sufficient conditions under which diminishing returns and competitive incentives generate catch-up. 
The key additional requirement is that a developer that trails in a capability retains marginal capacity to improve that capability at a rate related to the size of its gap.
Intuitively, lagging capabilities retain greater marginal capacity for improvement, while competition creates incentives for developers to perform well in lagging capabilities.

\begin{restatable}[Sufficient conditions for capability catch-up]{proposition}{PrimitiveCatchup}
\label{prop:PrimitiveCatchup}
Suppose each $F_i$ is continuously differentiable and coordinatewise
nondecreasing, each coordinate function $[F_i(\cdot)]_k$ is concave, and
$F_i$ is $L$-Lipschitz with respect to $\ell_1$.
Additionally, suppose $w_k>0$ with $w_{\mathrm{min}} = \min\limits_{k \in [K]} w_k$, and $c$ is twice continuously differentiable with $c'(0)=0$ and $0\le c''(s)\le M$. 
Assume that each developer is playing a pure-strategy Nash equilibrium $\mathbf{x}^*_{i,t}$ at each timestep.  

For each developer $i$, category $k$, and investment state $\mathbf x$, define
\[
    m_{i,k}(\mathbf x)
    :=
    \min_{\ell\in[K]}
    \frac{\partial [F_i(\mathbf x)]_k}{\partial x_\ell}.
\]
Suppose there exists $\mu>0$ such that, for every pair of developers
$i,i'$, category $k$, and relevant investment states $\mathbf x,\mathbf y$,
whenever $[F_i(\mathbf x)]_k \le [F_{i'}(\mathbf y)]_k$, we have
\[ 
    m_{i,k}(\mathbf x)^2
    \ge
    \mu\left(
        [F_{i'}(\mathbf y)]_k-[F_i(\mathbf x)]_k
    \right).
\] 
Define
\[
    \underline p
    :=
    w_{\mathrm{min}}
    \min_{|z|\le \frac{L^2}{\mu}}\sigma'(z)>0.
\]
Then, for every developer $i$, category $k$, and period $t$,
\[
    \Delta\theta_{i,k}^t
    \ge
    \beta g_{i,k}^t,
    \qquad
    \beta
    :=
    \frac{\underline p\,\mu}
         {M+\underline p\,\mu}
    >0.
\] 
\end{restatable}

Proposition~\ref{prop:PrimitiveCatchup} illustrates how diminishing returns and competitive incentives can jointly generate catch-up. 
The parameter \(\mu\) captures how much marginal production capacity a developer retains relative to its capability gap, while \(\underline{p}\) lower-bounds the competitive value of improving a capability. 
As a result, the catch-up guarantee is stronger when lagging capabilities retain greater marginal production capacity, competitive evaluation places greater marginal value on improvement, and marginal investment costs grow more slowly. 
The diminishing returns assumption is plausible when capability gaps are bounded and lagging developers have not exhausted the channels available for improving that capability; it becomes restrictive when some investment inputs have little or no effect on a capability or when large gaps coexist with very small marginal returns. 
Diminishing returns allows lagging capabilities greater remaining capacity for improvement, while competition gives developers an incentive to exploit this capacity. 

The proof combines the lower bound on marginal production capacity with the developers' first-order conditions to obtain a lower bound on the capability improvement of a trailing developer (Appendix~\ref{app:PrimitiveCatchup}).
Because the cost function aggregates investment across dimensions, the equilibrium conditions determine the total amount of investment but not, absent solving the equilibrium, how that investment is distributed across inputs. 
A uniform catch-up guarantee therefore requires a lower bound that holds for every possible active investment direction, which leads naturally to the minimum marginal effect across dimensions.
Relaxing this assumption therefore requires a different cost function that places more structure on how investment is allocated across dimensions, or additional assumptions about which investment dimensions are active at equilibrium. 
As a result, this condition should be interpreted as a sufficient primitive condition for catch-up, rather than as a necessary condition for catch-up.

Sections~\ref{sec:monoculture} and~\ref{sec:dynamics} characterize two dimensions of convergence. 
Developers may become similar in their capability profiles (monoculture), or in their overall performance levels (crowding). 
Our empirical analysis asks how these two forms of convergence evolve over a finite period, and whether frontier crowding can coexist with persistent differences in relative capability profiles. 